\section{Implementation}

This chapter describes the implementation of the proposed system design.

\subsection{Technology stack}

The system is implemented using the following technologies:
\begin{itemize}
  \item Python 3.13 for the core processing logic;
  \item PostgreSQL 17 for data storage;
  \item Docker for containerized deployment;
  \item GitHub Actions for continuous integration.
\end{itemize}

\subsection{Architecture overview}

The implemented system follows the hybrid architecture selected in Chapter~4.
The core module handles data processing and analysis, while extension modules provide specialized functionality.

% === EXAMPLE: TikZ deployment architecture with colors and shapes ===
\begin{figure}[H]
    \centering
    \begin{tikzpicture}[
        box/.style={rectangle, draw, minimum width=2.5cm, minimum height=0.8cm, align=center},
        cloud/.style={ellipse, draw, minimum width=2cm, minimum height=0.7cm, align=center},
        >=Stealth
    ]
        \node[cloud] (client) {Client};
        \node[box, right=1.5cm of client] (lb) {Load Balancer};
        \node[box, above right=0.5cm and 1.5cm of lb] (app1) {App Server 1};
        \node[box, below right=0.5cm and 1.5cm of lb] (app2) {App Server 2};
        \node[box, right=1.5cm of lb, yshift=0cm] (cache) {Cache};
        \node[box, right=3.5cm of lb] (db) {Database};
        \draw[->] (client) -- (lb);
        \draw[->] (lb) -- (app1);
        \draw[->] (lb) -- (app2);
        \draw[->] (app1) -- (cache);
        \draw[->] (app2) -- (cache);
        \draw[->] (cache) -- (db);
    \end{tikzpicture}
    \caption{Implemented system architecture} \label{fig:impl-arch}
\end{figure}

% === EXAMPLE: Including an external image with \addimghere ===
\cref{fig:sample-photo} demonstrates including an external image from the \texttt{images/} directory using the \verb|\addimghere| custom command.

\addimghere{sample-photo}{0.4}{Sample image included from file}{fig:sample-photo}

% === EXAMPLE: Two images side by side with \addtwoimghere ===
\cref{fig:two-images} shows two images placed side by side using the \verb|\addtwoimghere| command.
This is useful for comparing visual results or showing before/after states.

\addtwoimghere{sample-photo}{0.4}{sample-photo}{0.4}{Two images side by side}{fig:two-images}

\subsection{Key components}

The main processing module is responsible for data ingestion and transformation.
\cref{lst:processor} shows a simplified version of the data processor.

% === EXAMPLE: Python code listing with line numbers ===
\begin{lstlisting}[language=Python, caption={Data processor module}, label={lst:processor}]
import json
import logging
from dataclasses import dataclass
from typing import List, Optional

logger = logging.getLogger(__name__)

@dataclass
class DataPoint:
    '''Represents a single data measurement.'''
    timestamp: str
    value: float
    label: Optional[str] = None

class DataProcessor:
    '''Processes and transforms input data.'''

    def __init__(self, config: dict):
        self.config = config
        self.results: List[dict] = []

    def load_data(self, filepath: str) -> List[DataPoint]:
        '''Load data from a JSON file.'''
        with open(filepath, 'r') as f:
            raw = json.load(f)
        return [DataPoint(**item) for item in raw]

    def process(self, data: List[DataPoint]) -> List[dict]:
        '''Apply transformations to the data.'''
        for point in data:
            result = {
                'timestamp': point.timestamp,
                'processed_value': point.value * self.config.get('scale', 1.0),
                'label': point.label or 'unlabeled',
            }
            self.results.append(result)
        logger.info(f"Processed {len(self.results)} data points")
        return self.results

    def export(self, output_path: str) -> None:
        '''Export results to a JSON file.'''
        with open(output_path, 'w') as f:
            json.dump(self.results, f, indent=2)
        logger.info(f"Results exported to {output_path}")
\end{lstlisting}

\subsection{Configuration management}

The system uses a configuration file to manage runtime parameters.
\cref{lst:config} shows an example configuration.

% === EXAMPLE: JSON code listing ===
\begin{lstlisting}[caption={Example configuration file (config.json)}, label={lst:config}]
{
  "scale": 1.5,
  "output_dir": "./results",
  "log_level": "INFO",
  "database": {
    "host": "localhost",
    "port": 5432,
    "name": "thesis_db"
  },
  "processing": {
    "batch_size": 100,
    "max_retries": 3,
    "timeout_seconds": 30
  }
}
\end{lstlisting}

\subsection{Database schema}

% === EXAMPLE: SQL code listing ===
\cref{lst:sql} shows the database schema used for storing experimental results.

\begin{lstlisting}[language=SQL, caption={Database schema (schema.sql)}, label={lst:sql}]
CREATE TABLE experiments (
    id          SERIAL PRIMARY KEY,
    name        VARCHAR(255) NOT NULL,
    created_at  TIMESTAMP DEFAULT CURRENT_TIMESTAMP,
    status      VARCHAR(50) DEFAULT 'pending'
);

CREATE TABLE results (
    id            SERIAL PRIMARY KEY,
    experiment_id INTEGER REFERENCES experiments(id),
    metric_name   VARCHAR(100) NOT NULL,
    metric_value  DOUBLE PRECISION NOT NULL,
    recorded_at   TIMESTAMP DEFAULT CURRENT_TIMESTAMP
);

CREATE INDEX idx_results_experiment
    ON results(experiment_id);
\end{lstlisting}

\subsection{Dockerfile}

% === EXAMPLE: Dockerfile listing ===
\cref{lst:dockerfile} shows the Dockerfile used for containerized deployment.

\begin{lstlisting}[caption={Application Dockerfile}, label={lst:dockerfile}]
FROM python:3.13-slim
WORKDIR /app
COPY requirements.txt .
RUN pip install --no-cache-dir -r requirements.txt
COPY src/ src/
COPY config/ config/
EXPOSE 8080
CMD ["python", "-m", "src.main"]
\end{lstlisting}

\subsection{Security architecture}

The system implements several security measures, as shown in \cref{fig:security}.

% === EXAMPLE: TikZ diagram with colored fills ===
\begin{figure}[H]
    \centering
    \begin{tikzpicture}[
        box/.style={rectangle, draw, minimum width=2.5cm, minimum height=0.8cm, align=center},
        >=Stealth
    ]
        \node[box, fill=blue!10] (fw) {Firewall};
        \node[box, right=1cm of fw, fill=green!10] (auth) {Authentication};
        \node[box, right=1cm of auth, fill=orange!10] (enc) {Encryption};
        \node[box, below=0.8cm of auth, fill=yellow!10] (audit) {Audit Logging};
        \draw[->] (fw) -- (auth);
        \draw[->] (auth) -- (enc);
        \draw[->] (auth) -- (audit);
    \end{tikzpicture}
    \caption{Security architecture overview} \label{fig:security}
\end{figure}

\clearpage
